\textbf{Phần 1. Trắc nghiệm lựa chọn}
\setcounterpageref{ex}{0}

\begin{ex} %[1D8V32]
Tính kết quả của phép tính: $58 - 23$.
\choice
{$34$}
{\True $35$}
{$36$}
{$33$}
\loigiai{
Hàng đơn vị: $8 - 3 = 5$. Hàng chục: $5 - 2 = 3$. Vậy $58 - 23 = 35$.
\textbf{Đáp số:} $35$.\par
\textbf{Đáp án của câu hỏi là B}
}
\end{ex}

\begin{ex} %[1D8V32]
Tính kết quả của phép tính: $76 - 34$.
\choice
{$41$}
{$43$}
{$40$}
{\True $42$}
\loigiai{
Hàng đơn vị: $6 - 4 = 2$. Hàng chục: $7 - 3 = 4$. Vậy $76 - 34 = 42$.
\textbf{Đáp số:} $42$.\par
\textbf{Đáp án của câu hỏi là D}
}
\end{ex}

\begin{ex} %[1D8V32]
Tính kết quả của phép tính: $89 - 45$.
\choice
{\True $44$}
{$43$}
{$45$}
{$42$}
\loigiai{
Hàng đơn vị: $9 - 5 = 4$. Hàng chục: $8 - 4 = 4$. Vậy $89 - 45 = 44$.
\textbf{Đáp số:} $44$.\par
\textbf{Đáp án của câu hỏi là A}
}
\end{ex}

\begin{ex} %[1D8V32]
Điền số thích hợp: $67 - \dots = 32$.
\choice
{$34$}
{$36$}
{\True $35$}
{$33$}
\loigiai{
Số cần điền là: $67 - 32 = 35$.
\textbf{Đáp số:} $35$.\par
\textbf{Đáp án của câu hỏi là C}
}
\end{ex}

\begin{ex} %[1D8V32]
Điền số thích hợp: $\dots - 24 = 53$.
\choice
{$76$}
{\True $77$}
{$78$}
{$75$}
\loigiai{
Số cần điền là: $53 + 24 = 77$.
\textbf{Đáp số:} $77$.\par
\textbf{Đáp án của câu hỏi là B}
}
\end{ex}

\begin{ex} %[1D8V32]
Trường Tiểu học có 85 bạn học sinh khối 1. Trong đó có 42 bạn nam. Hỏi khối 1 có bao nhiêu bạn nữ?
\choice
{$42$}
{$44$}
{$41$}
{\True $43$}
\loigiai{
Số học sinh nữ khối 1 là: $85 - 42 = 43$ (bạn).
\textbf{Đáp số:} $43$.\par
\textbf{Đáp án của câu hỏi là D}
}
\end{ex}

\begin{ex} %[1D8V32]
Một cửa hàng có 96 chiếc quạt điện, đã bán được 51 chiếc quạt. Hỏi cửa hàng còn lại bao nhiêu chiếc quạt điện?
\choice
{\True $45$}
{$44$}
{$46$}
{$43$}
\loigiai{
Số quạt điện còn lại là: $96 - 51 = 45$ (chiếc quạt).
\textbf{Đáp số:} $45$.\par
\textbf{Đáp án của câu hỏi là A}
}
\end{ex}

\begin{ex} %[1D8V32]
Bác nông dân có 78 kg gạo, bác đã bán đi 35 kg gạo. Hỏi bác nông dân còn lại bao nhiêu ki-lô-gam gạo?
\choice
{$42$}
{$44$}
{\True $43$}
{$41$}
\loigiai{
Số ki-lô-gam gạo còn lại là: $78 - 35 = 43$ (kg).
\textbf{Đáp số:} $43$.\par
\textbf{Đáp án của câu hỏi là C}
}
\end{ex}

\begin{ex} %[1D8V32]
Một sợi dây thừng dài 88 cm. Người ta cắt đi một đoạn dài 36 cm. Hỏi sợi dây thừng còn lại dài bao nhiêu xăng-ti-mét?
\choice
{$51\text{ cm}$}
{\True $52\text{ cm}$}
{$53\text{ cm}$}
{$50\text{ cm}$}
\loigiai{
Chiều dài còn lại của sợi dây là: $88 - 36 = 52$ (cm).
\textbf{Đáp số:} $52$.\par
\textbf{Đáp án của câu hỏi là B}
}
\end{ex}

\begin{ex} %[1D8V32]
Tính giá trị của biểu thức: $98 - 32 - 24$.
\choice
{$41$}
{$43$}
{$40$}
{\True $42$}
\loigiai{
Thực hiện: $98 - 32 = 66$; $66 - 24 = 42$.
\textbf{Đáp số:} $42$.\par
\textbf{Đáp án của câu hỏi là D}
}
\end{ex}

\begin{ex} %[1D8V32]
Tính giá trị của biểu thức: $75 - 20 - 15$.
\choice
{\True $40$}
{$30$}
{$50$}
{$39$}
\loigiai{
Thực hiện: $75 - 20 = 55$; $55 - 15 = 40$.
\textbf{Đáp số:} $40$.\par
\textbf{Đáp án của câu hỏi là A}
}
\end{ex}

\begin{ex} %[1D8V32]
Tìm số tự nhiên $x$ thỏa mãn: $x + 31 = 86$.
\choice
{$54$}
{$56$}
{\True $55$}
{$53$}
\loigiai{
Ta có: $x = 86 - 31 = 55$.
\textbf{Đáp số:} $55$.\par
\textbf{Đáp án của câu hỏi là C}
}
\end{ex}

\textbf{Phần 2. Trắc nghiệm đúng sai}
\setcounterpageref{ex}{0}

\begin{ex} %[1D8V32]
Xét tính đúng hoặc sai của các khẳng định toán học sau đây:
\choiceTF
{Kết quả của bài toán «Tìm số tự nhiên $y$ thỏa mãn: $y + 43 = 97$» là $55$.}
{\True Đàn vịt có 65 con. Có 32 con vịt bơi sang ao bên cạnh. đàn vịt còn lại $33$ con vịt.}
{Một bãi đỗ xe có 59 chiếc ô tô, có 24 chiếc ô tô rời khỏi bãi. trong bãi còn lại $34$ chiếc ô tô.}
{\True Kết quả của bài toán «Hiệu giữa số lớn nhất có hai chữ số và số 34 bằng bao nhiêu» là $65$.}
\loigiai{
\begin{itemize}
    \item Ý a) Sai: Kết quả đúng là $54$, không phải $55$. Ta có: $y = 97 - 43 = 54$.
\textbf{Đáp số:} $54$.
    \item Ý b) Đúng: Số con vịt còn lại là: $65 - 32 = 33$ (con).
\textbf{Đáp số:} $33$.
    \item Ý c) Sai: Kết quả đúng là $35$, không phải $34$. Số ô tô còn lại là: $59 - 24 = 35$ (chiếc ô tô).
\textbf{Đáp số:} $35$.
    \item Ý d) Đúng: Số lớn nhất có hai chữ số là 99: $99 - 34 = 65$.
\textbf{Đáp số:} $65$.
\end{itemize}
\begin{itemize}
    \item a) \textbf{Sai}.
    \item b) \textbf{Đúng}.
    \item c) \textbf{Sai}.
    \item d) \textbf{Đúng}.
\end{itemize}
}
\end{ex}

\begin{ex} %[1D8V32]
Xét tính đúng hoặc sai của các bài toán và phép tính sau đây:
\choiceTF
{\True Kết quả của bài toán «Hiệu giữa số 87 và số tròn chục nhỏ nhất có hai chữ số là bao nhiêu» là $77$.}
{\True Một cuốn sách có 95 trang. Mai đã đọc được 43 trang. Mai còn $52$ trang chưa đọc.}
{Hai bạn An và Bình có tổng cộng 78 viên bi. Biết An có 36 viên bi. Bình có $43$ viên bi.}
{Một bãi đỗ xe có 59 chiếc ô tô, có 24 chiếc ô tô rời khỏi bãi. trong bãi còn lại $34$ chiếc ô tô.}
\loigiai{
\begin{itemize}
    \item Ý a) Đúng: Số tròn chục nhỏ nhất có hai chữ số là 10: $87 - 10 = 77$.
\textbf{Đáp số:} $77$.
    \item Ý b) Đúng: Số trang Mai chưa đọc là: $95 - 43 = 52$ (trang).
\textbf{Đáp số:} $52$.
    \item Ý c) Sai: Kết quả đúng là $42$, không phải $43$. Số viên bi của Bình là: $78 - 36 = 42$ (viên bi).
\textbf{Đáp số:} $42$.
    \item Ý d) Sai: Kết quả đúng là $35$, không phải $34$. Số ô tô còn lại là: $59 - 24 = 35$ (chiếc ô tô).
\textbf{Đáp số:} $35$.
\end{itemize}
\begin{itemize}
    \item a) \textbf{Đúng}.
    \item b) \textbf{Đúng}.
    \item c) \textbf{Sai}.
    \item d) \textbf{Sai}.
\end{itemize}
}
\end{ex}

\textbf{Phần 3. Trả lời ngắn}
\setcounterpageref{ex}{0}

\begin{ex} %[1D8V32]
Số bi của Bình nhiều hơn số bi của An ở bài toán trên là bao nhiêu viên bi?
\shortans{6}
\loigiai{
\textbf{Đáp số:} $6$\par
Số bi của Bình nhiều hơn của An là: $42 - 36 = 6$ (viên bi).
\textbf{Đáp số:} $6$.
}
\end{ex}

\begin{ex} %[1D8V32]
Tìm số thích hợp: $77 - 23 = 11 + \dots$.
\shortans{43}
\loigiai{
\textbf{Đáp số:} $43$\par
Vế trái: $77 - 23 = 54$. Số cần điền: $54 - 11 = 43$.
\textbf{Đáp số:} $43$.
}
\end{ex}

\begin{ex} %[1D8V32]
So sánh: $86 - 41$ và $97 - 52$. Kết quả của mỗi phép tính đều bằng bao nhiêu?
\shortans{45}
\loigiai{
\textbf{Đáp số:} $45$\par
Ta có: $86 - 41 = 45$ và $97 - 52 = 45$. Cả hai phép tính đều bằng $45$.
\textbf{Đáp số:} $45$.
}
\end{ex}

\begin{ex} %[1D8V32]
Một rổ cam có 68 quả. Mẹ bán đi 25 quả và biếu bà 12 quả. Hỏi trong rổ còn lại bao nhiêu quả cam?
\shortans{31}
\loigiai{
\textbf{Đáp số:} $31$\par
Tổng số cam đã bán và biếu là: $25 + 12 = 37$ (quả).\nSố cam còn lại là: $68 - 37 = 31$ (quả cam).
\textbf{Đáp số:} $31$.
}
\end{ex}

\begin{ex} %[1D8V32]
Tính: $69 - 14 - 23$.
\shortans{32}
\loigiai{
\textbf{Đáp số:} $32$\par
Thực hiện: $69 - 14 = 55$; $55 - 23 = 32$.
\textbf{Đáp số:} $32$.
}
\end{ex}

\begin{ex} %[1D8V32]
Một cây tre dài 98 cm, chặt bớt đi một đoạn dài 45 cm. Đoạn tre còn lại dài bao nhiêu xăng-ti-mét?
\shortans{53}
\loigiai{
\textbf{Đáp số:} $53$\par
Đoạn tre còn lại dài: $98 - 45 = 53$ (cm).
\textbf{Đáp số:} $53$.
}
\end{ex}
